\section{Problem Setting and Related Work}\label{sec:background}

\begin{table*}[t]
\centering
\begin{minipage}[t]{0.34\textwidth}
\centering
\small
\setlength{\tabcolsep}{3pt}
\begin{tabular}{@{}lcccc@{}}
\toprule
Method & Retro. & FT & Quant & Distill \\
\midrule
Crypto Hash & \cmark & \xmark & \xmark & \xmark \\
Metadata/Logs & \xmark & \cmark & \cmark & \cmark \\
Watermarking & \xmark & $\sim$ & $\sim$ & \xmark \\
Behavioral Fingerprints & \cmark & \cmark & \cmark & \xmark \\
\textbf{Ours} & \cmark & \cmark & \cmark & \cmark \\
\bottomrule
\end{tabular}
\captionof{table}{Comparison with prior provenance methods. Retro.: applicable retroactively. FT: survives fine-tuning. Quant: survives quantization. Distill: correctly classifies distilled copies as unrelated. $\sim$: partial or unreliable.}
\label{tab:comparison}
\end{minipage}
\hfill
\begin{minipage}[t]{0.62\textwidth}
\centering
\small
\begin{tabular}{@{}l@{\hspace{6pt}}p{3.8cm}@{\hspace{6pt}}l@{}}
\toprule
Architecture & Residual branch $F(x)$ & Product $M$ \\
\midrule
Transformer MLP / BasicBlock & $W_2 \sigma(W_1 x)$ & $W_2 W_1$ \\
Attention V/O path & $W_O \mathrm{Attn}(x) W_V x$ & $W_O W_V$ \\
Attention Q/K path & bilinear $x^\top W_Q W_K^\top x$ & $W_Q W_K^\top$ \\
Bottleneck ResNet & $W_3 \sigma(W_2 \sigma(W_1 x))$ & $W_3 W_2 W_1$ \\
SwiGLU MLP & $W_{\mathrm{down}} [\sigma(W_{\mathrm{gate}} x) \odot W_{\mathrm{up}} x]$ & $W_{\mathrm{down}} W_{\mathrm{up}}$ \\
\bottomrule
\end{tabular}
\captionof{table}{Architecture-aware residual branch products. Each $M$ composes the linear maps of one branch (dropping nonlinearities) so it maps the residual stream to itself. Factorization must match the architecture; incomplete products destroy the signal (Appendix~\ref{app:resnet-factorization}).}
\label{tab:factorization}
\end{minipage}
\end{table*}

\subsection{White-Box Model Lineage Verification}

A verifier receives two checkpoints $A$ and $B$ and determines whether they share weight-level lineage. We assume white-box access and compatible residual architectures (same $L$ and $d$); cross-architecture comparisons are out of scope. The verifier operates on weights alone: no training data, activations, or forward passes.

Checkpoints are \emph{related} if they share a common weight ancestor through fine-tuning, RLHF, pruning, quantization, or LoRA merging. They are \emph{unrelated} if initialized independently, even when architecture, data, task, or outputs match; a distilled student trained from scratch is unrelated to its teacher, even if it mimics the outputs perfectly. Linear merges induce partial lineage (Appendix~\ref{app:harder-bench}).\footnote{All appendix references refer to Technical Supplement.} The verifier returns \textsc{Related} or \textsc{Unrelated}. The score is symmetric, $\mathcal{L}(A,B) = \mathcal{L}(B,A)$; directional attribution requires metadata (Appendix~\ref{app:verification}). The method complements cryptographic signing and watermarking where those were never applied.

\subsection{Related Work}

Table~\ref{tab:comparison} summarizes the comparison.

Cryptographic hashes break after any weight modification. Model cards~\citep{mitchell2019model} require honest distributors and can be falsified when checkpoints are redistributed. Proof-of-learning~\citep{jia2021proofoflearning} requires retained training transcripts that already-released checkpoints lack.

Parameter watermarks~\citep{uchida2017embedding} and backdoor watermarks~\citep{adi2018turning} embed deliberate signals during training. \citet{lukas2022sok} found no surveyed scheme that reliably survived fine-tuning and pruning, and all require proactive insertion: unwatermarked legacy checkpoints cannot be verified.

IPGuard~\citep{cao2021ipguard} extracts inputs near classification boundaries; adversarial frontier stitching~\citep{lemerrer2020adversarial} marks decision surfaces with adversarial examples. CKA~\citep{kornblith2019similarity} and SVCCA~\citep{raghu2017svcca} compare representations on a probe dataset. All four measure functional similarity and cannot distinguish inheritance from imitation: a distilled student may share decision surfaces while having independent weights. They also require data and forward passes.

Weight cosine, Frobenius distance, and permutation matching~\citep{ainsworth2023git} detect derived checkpoints when the suspect remains close in parameter space, but these are uncalibrated global scalars with no structural account of lineage. HuRef~\citep{zeng2025huref} relies on LLM parameter convergence after pretraining. REEF~\citep{zhang2024reef} needs probe data or activations. MoTHer~\citep{horwitz2025mother} recovers model trees but requires a checkpoint collection rather than a single pair. Our fingerprint is data-free, operates on a single pair against a calibrated null, and localizes inheritance to specific blocks.

\citet{park2026idroppedneuralnet} observed that correctly paired projections in trained residual blocks produce trace-concentrated products and used this to recover projection correspondences. We build on that observation: rather than treating trace concentration as evidence of ancestry, we remove the identity-aligned component and compare checkpoint-specific signatures. This enables a model-level test of weight inheritance under post-training transformations, reparameterizations, partial inheritance, and distillation.
